\mwnewpart
\renewcommand\runningtitle{Section 8. BASIC}
\chapter{BASIC}
\sectioncontents

\section{Microsoft's GW-BASIC}\ixs{BASIC}\ixs{GW-BASIC}

\begin{plainwords}
IBM's BASIC for PC DOS 2.10 was a program on the diskette that used the
Cassette BASIC in the Personal Computer's read-only memory, and both were
IBM's. In 2020 Microsoft published the source of GW-BASIC, its BASIC for
other makers' computers, as it stood in 1983: the language of IBM's
\typed{BASICA}, with no need of a ROM. As with MS-DOS, Microsoft left out
the part each maker wrote for its own computer: here the screen,
keyboard, graphics, speaker, printers and serial ports. RetrOS/LT builds
Microsoft's part from the source, with the same assembler as the rest of
the system, and adds that part for the IBM Personal Computer.
\end{plainwords}

\begin{tabularx}{\textwidth}{@{}P{2.0in}Y@{}}
\toprule
\file{upstream/gwbasic} & Microsoft's repository, pinned as a submodule:
thirty-five modules and five headers\\
\file{patches/gwbasic} & Two patches (Appendix~B)\\
\file{basic/PCIO.ASM} & The IBM PC part: sixty-three routines\\
\file{basic/BASIC.ASM} & \file{BASIC.COM}, which runs \file{BASICA.EXE}\\
\file{basic/samples} & The six sample programs\\
\file{scripts/prepbasic.py} & Makes the build tree\\
\file{scripts/buildbasic.sh} & Assembles and links\\
\bottomrule
\end{tabularx}

\section{Building It}\ixs{buildbasic.sh}\ixs{prepbasic.py}

The source was written for MASM 1.x. MASM 1.10, which the MS-DOS 2.0
release includes, wrote object files for twenty of the modules, six of
them cut short, and then ran on without end in \file{GWEVAL.ASM}. MASM
5.10 assembles all of them after four small changes, none to code
(\patch{GW 0001}). The thirteen modules MASM 1.10 finished were compared
byte for byte with MASM 5.10's: twelve are the same, and in
\file{GIOCOM.ASM} one instruction that refers forward to a constant has
the longer of its two encodings, the same addition.

The one change that mattered most is in \file{BINTRP.H}. Each statement's
address is declared external the first time the dispatch tables name it,
and a macro marks it as done. That macro was empty, closed with
\code{\&endm} so that the macro defining it would not end there; MASM 5
does not take \code{\&endm} as the end of anything, and the rest of every
file including \file{BINTRP.H} became part of the macro. Those modules
assembled with no errors and no code. An equate marks the name instead:

\srclist{gwb_bintrp}{BINTRP.H, as patched.}

\file{scripts/prepbasic.py} cleans the files as \file{prep.py} cleans the
MS-DOS source, joins \file{MATH1.ASM} and \file{MATH2.ASM} (one module,
split in two files), applies the patches, adds \file{PCIO.ASM} and
\file{PCIOID.INC}, which names the edition in the heading, and writes
everything with CR LF. \file{scripts/buildbasic.sh} assembles each module
and links them. The order matters three ways:

\srclist{buildbasic_order}{The link order.}

\engnote{\code{LSTVAR}, the last byte of the data segment, is where
BASIC's program text begins, and the byte before it must stay 0: it is
the end of the ``line'' before the first, and \typed{RUN} and
\typed{GOTO} start reading there. The first build put \file{PCIO.ASM}'s
data last, and its last byte was a counter. Every \typed{RUN} after the
counter first moved said \shown{Syntax error}, with no line number. A
byte of 0 now ends \file{PCIO.ASM}'s data:}

\srclist{pcio_guard}{The end of PCIO.ASM's data.}

\section{The IBM PC Part}\ixs{PCIO.ASM}

Microsoft's modules call sixty-three routines they do not define. The
interface is not documented anywhere but in the comments at the calls,
which say what each routine is given and must return; \file{PCIO.ASM}
follows them. Its routines fall into seven groups:

\begin{tabularx}{\textwidth}{@{}P{1.0in}Y@{}}
\toprule
Screen & \routine{SCROUT}, \routine{SCRINP}, \routine{SCROLL},
\routine{CLRSCN}, \routine{CLREOL}, \routine{CSRDSP}, \routine{SCRSTT}
(\typed{SCREEN}), \routine{SWIDTH} (\typed{WIDTH}), \routine{SETCLR}
(\typed{COLOR}) and the function key line, through interrupt 10h.\\
Keyboard & \routine{KEYINP}, which turns the ROM's keys into
Microsoft's codes, \routine{INKMAP} and \routine{INFMAP}, which turn them
back into IBM's for \typed{INKEY\$} and \typed{INPUT\$}, and
\routine{MAPSUP}, for \keys{Alt} and a letter.\\
Graphics & The current point and its moves, \routine{SETC},
\routine{READC}, \routine{NSETCX}, PAINT's \routine{SCANR} and
\routine{SCANL}, GET and PUT, in the Color/Graphics Adapter's memory.\\
Sound & \routine{DONOTE}: a queue of notes, played by the timer.\\
Printers, serial & \routine{SNDLPT}, through interrupt 17h; the serial
ports' own chips, with a receive buffer for each.\\
Pen, joysticks & Interrupt 10h function 4; port 201h.\\
Start, end & \routine{GWINI}, \routine{GWTERM}, \routine{SEGINI},
\routine{SETCBF}, \routine{GETHED}.\\
\bottomrule
\end{tabularx}

\subsection{The keyboard} BASIC's editor works in Microsoft's codes: an
editing key is two bytes, 0FFh and a control character (cursor up is
0FFh, 30), a function key 80h and 20h plus its number. \routine{KEYINP}
makes them from the ROM's scan codes:

\srclist{pcio_edkeys}{The editing keys.}

\engnote{\key{Esc} becomes 0FFh, 21, the editor's ``erase the line'': the
screen editor treats the character 27 as the start of an escape sequence
before it asks what the key is for. \routine{INKMAP} gives
\typed{INKEY\$} the 27 back.}

\subsection{PAINT} PAINT is Microsoft's, a scan-line fill with a queue;
the machine's part scans one line. \routine{SCANR} is given how many
points of the line above (or below) remain to be covered, skips border
points among them, then paints to the next border:

\srclist{pcio_scanr}{SCANR.}

\subsection{Sound} \typed{SOUND}, \typed{PLAY} and \typed{BEEP} give
\routine{DONOTE} a frequency and a duration in units of 2.5 ms.
\routine{DONOTE} queues the note (32 at most), and, if the timer is at its
usual 18.2 interrupts a second, sets it to 400: one interrupt a unit. The
BIOS's own routine still gets its 18.2 a second:

\srclist{pcio_tmrisr}{The timer interrupt, while music plays.}

When the queue is empty and the last note is over, the timer goes back to
18.2. In the foreground (\typed{MF}), \routine{DONOTE} waits for the note
to finish, or for \keys{Ctrl}{Break}.

\section{BASIC /ROM}\ixs{Cassette BASIC}\ixs{BASIC /ROM}

\routine{GWINI} looks for \typed{/ROM} in the command line before
anything else is done, so nothing has to be put back. The IBM Personal
Computer, XT, AT and PS/2 have Cassette BASIC at F600:0000 and
interrupt 18h pointing there; a computer without it points interrupt 18h
at its own routine (a message, or another try at starting):

\srclist{pcio_romchk}{/ROM.}

\engnote{The check of the first byte was added after a run on MAME's
5150, whose BASIC sockets are filled with zeros for the tests: interrupt
18h pointed at F600:0000, and the first version jumped into the zeros.
Empty sockets on a real Personal Computer read FFh. A checksum of the
ROM would be surer on the 5150, where the power-on test keeps each 8 KB
summing to 0, but the PS/2's BASIC is part of a larger ROM and need not.}

Cassette BASIC takes the whole computer, and DOS is gone; it has no
disk, only the cassette port, which only the original Personal Computer
has. The \textit{Guide to Operations} says so where it describes the
switch.

\section{BASIC.COM}\ixs{BASIC.COM}

PC DOS 2.10 had \file{BASIC.COM} and \file{BASICA.COM}. RetrOS/LT has one
BASIC, \file{BASICA.EXE} (an \file{.EXE}: Microsoft's \file{BIBOOT.ASM}
says it must not be converted to a \file{.COM}), and a 544-byte
\file{BASIC.COM} that runs it with the same command line. DOS 2 does not
tell a program where it was loaded from, and its EXEC function does not
search the path, so \file{BASIC.COM} tries the current directory and
then each directory of \code{PATH}, as \file{COMMAND.COM} does.

\section{What Was Added, and What Was Not}

\begin{itemize}
\item \textbf{PLAY's \code{>} and \code{<}}, the octave commands of IBM
  BASIC 2.0, were not in the 1983 source; \patch{GW 0002} adds them.
\item \textbf{The statements of later versions}, \typed{TIMER},
  \typed{CHDIR}, \typed{MKDIR}, \typed{RMDIR}, \typed{SHELL},
  \typed{ENVIRON}, \typed{IOCTL}, \typed{ERDEV}, \typed{VIEW},
  \typed{WINDOW} and \typed{PMAP}, are in the source as names only and
  answer \shown{Advanced Feature}; PC DOS 2.10's \typed{BASICA} did not
  have them either. The function \typed{PLAY(}\var{n}\typed{)} is not in
  the source.
\item \textbf{The cassette device} is left out by the source's own switch
  (\code{GWCASS=0}).
\item \textbf{Sequential files} are written in whole 128-byte records,
  the last ending in \^{}Z, as this BASIC always did.
\end{itemize}
